\section{LP spectral scales and the Newton right-hand side}
\label{sec:lp-spectra}

The diameter law determines the ratio of the extreme eigenvalues. For
a polytope, the dimension of the optimal face also determines how many
eigenvalues have each of two orders. For the logarithmic barrier this
splitting is classical active-set geometry, closely related to the
Hessian analyses of \citet{wright1994} and \citet{forsgren2002}. Using
\cref{thm:difference-body}, we show that every fixed finite-parameter
barrier has the same ordered eigenvalue orders. The eigenspaces need not
coincide, and their different rates of convergence matter for the
Newton equation.

\subsection{The optimal-face tangent and the two spectral scales}

We use the compact, strictly feasible standard-form representation
\eqref{eq:lp-standard}, now allowing an optimal face $\Phi$ of any
dimension $f$. Let $d=n-m$, and identify $\V=\ker A$ with $\R^d$ using
the orthonormal matrix $W$. Set
\[
 B=\{i:x_i>0\text{ for some }x\in\Phi\},\qquad
 N=\{1,\ldots,n\}\setminus B,\qquad
 \mathcal T=\{z\in\ker A:z_N=0\}.
\]
Then $\mathcal T$ is exactly the tangent space of $\aff\Phi$, and
$\dim\mathcal T=f$. Indeed, a relative-interior point $x^a$ of $\Phi$
is positive on $B$. Hence, for every sufficiently short
$z$ supported on $B$ with $Az=0$, both $x^a\pm z$ are
feasible. Optimality then forces the objective increment to vanish, so
both points lie in $\Phi$. Conversely, every chord of $\Phi$ vanishes
on $N$. The objective is orthogonal to this tangent:
$c_{\V}\perp\mathcal T$. Since the objective is nonconstant, $f<d$.
Write $P_{\mathcal T}$ for the orthogonal projector onto $\mathcal T$.

Classical endpoint theory for the logarithmic central path
\citep{adler1991,guler1994,halicka1999} gives
\begin{equation}
 x_B(\mu)\longrightarrow x_B^a>0,\qquad
 s_N(\mu)\longrightarrow s_N^a>0,
 \label{eq:lp-face-endpoints}
\end{equation}
where the limiting pair is strictly complementary, even when the LP is
degenerate. Positivity is componentwise. The reduced log-barrier Hessian
has the exact decomposition
\begin{equation}
 \begin{aligned}
 H_{\log}(x(\mu))&=H_N(\mu)+H_B(\mu),\\
 H_N&=\mu^{-2}W^\top R_N^\top\Diag(s_N^2)R_NW,\\
 H_B&=W^\top R_B^\top\Diag(x_B^{-2})R_BW,
 \end{aligned}
 \label{eq:lp-hard-soft}
\end{equation}
with coordinate restrictions $R_B,R_N$. Thus $\ker H_N=\mathcal T$
under the stated identification, while $\norm{H_B}=O(1)$.

\begin{theorem}[Two scales for every fixed barrier]
\label{thm:lp-all-spectra}
Let $F$ be any fixed $\nu$-self-concordant barrier on $P^\circ$.
In increasing eigenvalue order, as $g\downarrow0$,
\begin{equation}
 \lambda_j(\widehat H_F(g))=
 \begin{cases}
  \Theta(1),&1\leq j\leq f,\\
  \Theta(g^{-2}),&f<j\leq d.
 \end{cases}
 \label{eq:lp-all-eigenvalues}
\end{equation}
The first case is empty when the optimum is unique. All constants may
depend on the fixed instance and $\nu$. For sufficiently small $g$, let
$\Pi_F(g)$ be the orthogonal projector onto the eigenspace of the
lowest $f$ eigenvalues (the weak eigenspace); put $\Pi_F=0$ if $f=0$. Then
\begin{align}
 \norm{\Pi_F(g)-P_{\mathcal T}}&=O(g),
   \label{eq:general-projector-rate}\\
 \norm{\Pi_{\log}(g)-P_{\mathcal T}}&=O(g^2).
   \label{eq:log-projector-rate}
\end{align}
In particular, the weak eigenspace converges to $\mathcal T$,
although it need not equal $\mathcal T$ at positive gap.
\end{theorem}

\begin{proof}
The nonzero singular values of the fixed restriction $R_NW$ are
bounded away from zero. By \eqref{eq:lp-face-endpoints}, the restriction
of $H_N$ to $\mathcal T^\perp$ has all its eigenvalues between
$a\mu^{-2}$ and $b\mu^{-2}$, for fixed $a,b>0$ and all sufficiently small $\mu$.
Also $0\preceq H_B\preceq b'I$. Weyl's inequalities give $d-f$
eigenvalues of order $\mu^{-2}$. The minimax principle applied to
$\mathcal T$ bounds the lowest $f$ eigenvalues by $b'$. Dikin
containment in the compact set $P$ gives the common lower bound
$H_{\log}\succeq (\diam P)^{-2}I$. Therefore those $f$ eigenvalues
are of order one. By \eqref{eq:gap-mu-comparison}, $g=\Theta(\mu)$.

At equal objective gaps, \eqref{eq:equal-gap-loewner} gives constants
$a_F,b_F>0$ such that
\begin{equation}
 a_F\widehat H_{\log}(g)\preceq\widehat H_F(g)
 \preceq b_F\widehat H_{\log}(g).
 \label{eq:lp-reference-comparison}
\end{equation}
The eigenvalue minimax principle proves
\eqref{eq:lp-all-eigenvalues}. Moreover, for another constant $a'>0$,
\begin{equation}
 \widehat H_F(g)\succeq a'g^{-2}P_{\mathcal T^\perp}.
 \label{eq:hard-projector-lower}
\end{equation}
If $Q_F$ is an orthonormal basis for the weak eigenspace, its compressed
Hessian has norm $O(1)$. Evaluating \eqref{eq:hard-projector-lower}
on $Q_Fu$ gives
$\norm{P_{\mathcal T^\perp}Q_Fu}\leq O(g)\norm u$.

For the logarithmic barrier the eigenvector equation gives a sharper
estimate. Write $H_{\log}Q_{\log}=Q_{\log}\Lambda_w$, with
$\norm{\Lambda_w}=O(1)$. Then
\[
 H_NQ_{\log}=Q_{\log}\Lambda_w-H_BQ_{\log},
 \qquad \norm{H_NQ_{\log}}=O(1).
\]
The inverse of $H_N$ on $\mathcal T^\perp$ has norm $O(g^2)$, so
$\norm{P_{\mathcal T^\perp}Q_{\log}}=O(g^2)$.
For subspaces of the same dimension, the norm of the difference of
their orthogonal projectors equals the largest sine of their principal
angles, here $\norm{P_{\mathcal T^\perp}Q_F}$. This identity follows
by taking a singular value decomposition of the matrix of inner
products of their orthonormal bases. It proves both asserted projector
rates. This is the elementary spectral-subspace perturbation argument
underlying the familiar Davis--Kahan estimates \citep{daviskahan1970}.
\end{proof}

The proof uses endpoint convergence only for the reference logarithmic
barrier, not for the general barrier. By \cref{thm:difference-body}, the constants are
uniform over barriers with uniformly bounded $\nu$. The ordered
eigenvalue scales and the general $O(g)$ projector estimate also hold
at points whose Newton decrement \eqref{eq:centrality-residual} is at
most a fixed $\rho<1$ for some $\mu>0$, with $g$ the point's own
objective gap. The
right-hand-side statement below also requires exact centrality.

\subsection{The exact Newton right-hand side and discarding the weak eigenspace}

At an exact center $x_F(\mu_0)$, the Newton equation for
$F(x)+\inner c x/\mu_1$, with new parameter $\mu_1$, is
\begin{equation}
 H_F(x_F(\mu_0))\,\Delta x
 =-\bigl(\mu_1^{-1}-\mu_0^{-1}\bigr)c_{\V}.
 \label{eq:exact-newton-rhs}
\end{equation}
The path derivative has a right-hand side in the same direction,
since $H_F(x_F(\mu))x_F'(\mu)=\mu^{-2}c_{\V}$. Thus, for every
nonzero right-hand side $r$ in \eqref{eq:exact-newton-rhs},
\begin{equation}
 \frac{\norm{\Pi_F(g)r}}{\norm r}=O(g),\qquad
 \frac{\norm{\Pi_{\log}(g)r}}{\norm r}=O(g^2).
 \label{eq:newton-weak-fraction}
\end{equation}
These follow immediately from $P_{\mathcal T}c_{\V}=0$ and
\cref{thm:lp-all-spectra}. The quantities in
\eqref{eq:newton-weak-fraction} are ratios of \emph{norms}; the
corresponding ratios of squared norms are $O(g^2)$ and $O(g^4)$.
At an approximately central point, the right-hand side has an
additional centering residual, and spectral clustering alone does not
give these estimates for it.

For positive-definite $H$ and an orthogonal spectral projector
$\Pi$ of $H$, discarding the range of $\Pi$ gives
\begin{equation}
 u_s=H^{-1}(I-\Pi)r,\qquad
 r-Hu_s=\Pi r,\qquad
 H^{-1}r-u_s=H^{-1}\Pi r.
 \label{eq:discarding-identity}
\end{equation}
The relative Euclidean residual is thus exactly
$\norm{\Pi r}/\norm r$, but the relative solution error is not bounded
by this fraction. Now suppose, in the LP setting, that the spectral
projection and a representation of $H$ restricted to the complementary
eigenspace $\operatorname{range}(I-\Pi)$ of the large eigenvalues are
available. By
\cref{thm:lp-all-spectra}, this restricted operator has a bounded
condition number. If a computed restricted solution satisfies
\[
 \widetilde u_s\in\operatorname{range}(I-\Pi),\qquad
 \norm{(I-\Pi)r-H\widetilde u_s}
 \leq\delta\norm{(I-\Pi)r},
\]
then orthogonality gives
\begin{equation}
 \frac{\norm{r-H\widetilde u_s}^2}{\norm r^2}
 \leq\frac{\norm{\Pi r}^2}{\norm r^2}+\delta^2.
 \label{eq:conditional-window-residual}
\end{equation}
This bound ignores the cost of constructing the projection and does
not show that a path-following algorithm's own error criterion is
met. A Euclidean residual bound alone implies neither accurate Newton
directions nor convergence of an inexact interior-point method.

\subsection{A sharp general-barrier example}

The map $(x,y)\mapsto(1+x,1-x,y,1-y)$ embeds the rectangle below in
a strictly feasible standard form. It multiplies the tangent metric by
two, so it preserves eigenvector angles and all gap exponents. The
example therefore tests the rates of \cref{thm:lp-all-spectra}
directly.

Oscillating perturbations of self-concordant barriers are not new:
\citet[Remark 2.3]{duistermaat2001} uses $f+A\sin f$ to produce
rescaled derivatives without boundary limits. We use the perturbation
$x\sin(\log y)$ below, which couples the two coordinates, to show that
the eigenspace rate is sharp and to compare the residual with the
solution error.

The following finite-parameter barrier shows that the $O(g)$ norm rate
cannot in general be improved to $O(g^2)$. It also shows that the
small residual in \eqref{eq:discarding-identity} can come with a large
solution error.

\begin{proposition}[Sharp eigenspace rate, and a small residual with a large solution error]
\label{prop:oscillatory-barrier}
On $P=[-1,1]\times[0,1]$, minimize $y$ and set
\begin{equation}
 \begin{aligned}
 F(x,y)&=4\bigl(F_0(x,y)+\varepsilon\psi(x,y)\bigr),\\
 F_0(x,y)&=-\log(1-x^2)-\log y-\log(1-y),\\
 \psi(x,y)&=x\sin(\log y),\qquad\varepsilon=1/100.
 \end{aligned}
 \label{eq:oscillatory-barrier}
\end{equation}
This is a $20$-self-concordant barrier. At gaps
$g_k=\exp(-2\pi k)$, $k=1,2,\ldots$, its exact center is $(0,g_k)$.
For its weak spectral projector and $r=(0,1)^\top$,
\begin{equation}
 \norm{\Pi_F(g_k)r}\sim\varepsilon g_k.
 \label{eq:oscillatory-sharpness}
\end{equation}
With $H=\widehat H_F(g_k)$ and $u_s=H^{-1}(I-\Pi_F(g_k))r$,
\begin{equation}
 \frac{\norm{r-Hu_s}}{\norm r}\sim\varepsilon g_k,
 \qquad
 \frac{\norm{H^{-1}r-u_s}}{\norm{H^{-1}r}}\longrightarrow1.
 \label{eq:residual-direction-separation}
\end{equation}
\end{proposition}

\begin{proof}
We first check the barrier assumptions on the entire open rectangle.
The four affine logarithms comprising $F_0$ make it a standard
self-concordant barrier with parameter at most $4$. Its Hessian is
\[
 H_0=\Diag(A_0,B_0),\qquad
 A_0=\frac{2(1+x^2)}{(1-x^2)^2}\geq2,\qquad
 B_0=y^{-2}+(1-y)^{-2}\geq y^{-2}.
\]
For $h=(u,v)$ put $U=\sqrt{A_0}u$, $V=\sqrt{B_0}v$, so that
$\norm h_{H_0}^2=U^2+V^2$, $|u|\leq |U|/\sqrt2$, and
$|v|/y\leq |V|$. Writing $s=\sin(\log y)$ and
$t=\cos(\log y)$, direct differentiation gives
\begin{align*}
 D\psi[h]&=su+xtv/y,\\
 D^2\psi[h,h]&=2tuv/y-x(s+t)v^2/y^2,\\
 D^3\psi[h,h,h]&=-3(s+t)uv^2/y^2+x(3s+t)v^3/y^3.
\end{align*}
Since $|x|<1$, $|s+t|\leq\sqrt2$, and
$|3s+t|\leq\sqrt{10}$, these imply
\begin{align}
 |D\psi[h]|&\leq2\norm h_{H_0},\nonumber\\
 |D^2\psi[h,h]|&\leq3\norm h_{H_0}^2,\label{eq:perturbation-bounds}\\
 |D^3\psi[h,h,h]|&\leq7\norm h_{H_0}^3.\nonumber
\end{align}
For example, the second bound follows from
$\sqrt2|UV|+\sqrt2 V^2\leq3(U^2+V^2)$; the third follows from
$3|U|V^2+\sqrt{10}|V|^3\leq7(U^2+V^2)^{3/2}$.
For $f_0=F_0+\varepsilon\psi$ we therefore have
\[
 D^2f_0[h,h]\geq(1-3\varepsilon)\norm h_{H_0}^2,
 \quad |D^3f_0[h,h,h]|\leq(2+7\varepsilon)\norm h_{H_0}^3,
 \quad |Df_0[h]|\leq(2+2\varepsilon)\norm h_{H_0}.
\]
The inequality
$2+7\varepsilon\leq4(1-3\varepsilon)^{3/2}$ proves standard
self-concordance of $F=4f_0$. Its barrier parameter is at most
\begin{equation}
 \frac{4(2+2\varepsilon)^2}{1-3\varepsilon}<20.
 \label{eq:oscillatory-parameter}
\end{equation}
Its Hessian is positive definite. The bounded perturbation $\psi$
does not affect divergence at the boundary, completing the global check.

At $y=g_k$, $s=0$ and $t=1$. At $x=0$ the first derivative in $x$
vanishes, while
\[
 F_y(0,g_k)=4\left(-\frac1{g_k}+\frac1{1-g_k}\right)<0.
\]
Consequently $\mu_k=-1/F_y(0,g_k)>0$ makes the point stationary
for $F+y/\mu_k$. Strict convexity makes it the unique exact center.
The objective gap is $g_k$. Define $J=H/4$ to keep the scaling explicit:
\begin{equation}
 J=\begin{pmatrix}2&\varepsilon/g_k\\
             \varepsilon/g_k&g_k^{-2}+(1-g_k)^{-2}\end{pmatrix}.
 \label{eq:oscillatory-hessian}
\end{equation}
For $g=g_k\downarrow0$, its eigenvalues satisfy
$\lambda_+(J)\sim g^{-2}$ and
$\lambda_-(J)\to2-\varepsilon^2$, as follows from its trace and
determinant. A unit weak eigenvector $v_-$ with positive first
coordinate obeys
\[
 \frac{(v_-)_2}{(v_-)_1}
 =-\frac{2-\lambda_-(J)}{\varepsilon/g}
 \sim-\varepsilon g.
\]
It follows that $|v_-^\top r|\sim\varepsilon g$, whereas
$|v_+^\top r|\to1$. Restoring $H=4J$ gives
\begin{equation}
 \norm{H^{-1}\Pi_F r}\sim
   \frac{\varepsilon g}{4(2-\varepsilon^2)},\qquad
 \norm{H^{-1}(I-\Pi_F)r}\sim\frac{g^2}{4}.
 \label{eq:oscillatory-solution-components}
\end{equation}
The weak component thus dominates the Euclidean solution norm, while
the residual of $u_s$ is exactly $\Pi_Fr$.
This proves \eqref{eq:residual-direction-separation}.
\end{proof}

The example also shows why \cref{thm:lp-all-spectra} does not assume
that a general central path converges. At every
sufficiently small $y$, the unique solution of
\[
 \frac{2x}{1-x^2}+\varepsilon\sin(\log y)=0
\]
has $F_y(x,y)<0$ and hence lies on the exact path for a positive $\mu$.
Indeed, $F_y/4\leq(-1+\varepsilon)/y+1/(1-y)<0$ for small $y$.
The solution $x$ takes distinct limiting values along sequences with
$\sin(\log y)=1$ and $-1$. Thus the path has no single endpoint,
although the barrier parameter is finite. The weak eigenspace still
converges to $\mathcal T$, as the theorem asserts.

Nonconvergent central paths were already constructed by
\citet{boltepauwels2022} using a Legendre function that is finite and
continuous on a closed square. That function does not tend to
infinity at the boundary, as \cref{def:barrier} requires. Our example
gives nonconvergence for a self-concordant barrier whose parameter
is bounded on the whole rectangle by \eqref{eq:oscillatory-parameter},
together with the sharp projector and solution-error estimates.
